\begin{enunciado}{\ejercicio}
  Factorizar en
  $\complejos[X], \reales[X], \racionales[X]$ y en $(\enteros/7\enteros)[X]$ los polinomios cuadráticos
  \begin{enumerate}[label=\roman*)]
    \begin{multicols}{3}
      \item $X^2 + 6X - 2$,
      \item $X^2 + X - 6$,
      \item $X^2 - 2X + 10$
    \end{multicols}
  \end{enumerate}
\end{enunciado}

Fácil, si tenemos un polinomio de grado 2 de la pinta $aX^2 + bX + c$, mandamos la resolvente siempre con el cuidado de no escribir algo negativo dentro
de la raíz, es decir las raíces $r$:
$$
  r = \frac{-b \pm \omega}{2a} ~~ \text{con} ~~ \omega^2 = b^2 - 4ac.
$$
\parrafoDestacado{
  \tiny \it
  Cuenta la leyenda que si ponés un número negativo en el radicando, viene el Ratón Pérez, te caga a palos, te roba los dientes y se los vende a Satanas.
}

\begin{enumerate}[label=\roman*)]
  \item
        Haciendo las cuentas que no pienso hacer, las raíces:
        $$
          \omega = 2 \sqrt{11} \entonces r_{1,2} = (-3 \pm \sqrt{11}) \notin \racionales
        $$

        \textit{Con cuerpo $\K = \racionales$}:

        Con el resultado de la resolvente concluyo que el polinomio $X^2 + 6X - 2$ es \textit{irreducible} en $\racionales$.

        \bigskip

        \textit{Con cuerpo $\K = \reales, \complejos$}:

        Con el resultado de la resolvente:
        $$
          \cajaResultado{
            X^2 + 6X - 2 = (X - (-3 + \sqrt{11}) ) \cdot (X - (-3 - \sqrt{11}) )
          }
        $$

        \bigskip

        \textit{Con cuerpo $\K = (\enteros/7\enteros)$}:

        El polinomio $X^2 + 6X - 2$ en este cuerpo lo puedo escribir así:
        $$
          X^2 + 6X - 2  \Entonces{$\K = (\enteros/7\enteros)$} X^2 + 6X + 5
        $$
        Calcular las raíces de eso:
        $$
          r_{1,2} = -3 \pm 2 \conga7
          \llave{rcl}{
            r_1 & = & 6\\
            r_2 & = & 2
          }
        $$
        \parrafoDestacado{
          \it
          Fijate que si agarrabas las raíces previamente calculadas:
          $$
            -3 \pm \sqrt{11} \conga7 4 \pm \sqrt{4} = 4 \pm 2 =
            \llave{rcl}{
              r_1 & = & 6\\
              r_2 & = & 2
            }
          $$
        }
        La factorización buscada en $(\enteros/7\enteros)$:
        $$
          X^2 + 6X - 2 = (X - 6) (X - 2)
        $$

  \item
        Haciendo las cuentas que no pienso hacer, las raíces:
        $$
          \omega = 5 \entonces r_{1,2} =
          \llave{rcl}{
            r_1 & = & 2\\
            r_2 & = & -3
          }
        $$
        Resultados más lindo no podemos pedir.

        \textit{Con cuerpo $\K = \racionales, \reales, \complejos$}
        $$
          \cajaResultado{
            X^2 + X - 6 = (X -2) (X + 3)
          }
        $$

        \bigskip

        \textit{Con cuerpo $\K = (\enteros/7\enteros)$}:
        $$
          \cajaResultado{
            X^2 + X - 6 =
            X^2 + X + 1 =
            (X + 5)(X + 3)
          }
        $$

  \item
        Haciendo las cuentas que no pienso hacer, las raíces:
        $$
          \omega^2 = -36 \entonces r_{1,2} =
          \llave{rcl}{
            r_1 & = & 1 + 3i\\
            r_2 & = & 1 - 3i
          }
        $$

        \bigskip

        \textit{Con cuerpo $\K = \racionales, \reales$}:

        El polinomio $X^2 + -2X + 10$ es irreducible.

        \bigskip

        \textit{Con cuerpo $\K = \complejos$}:

        El polinomio tiene una factorización:
        $$
          \cajaResultado{
            X^2 -2X + 10 = (X - (1+3i))(X - (1 - (1-3i)))
          }
        $$

        \bigskip

        \textit{Con cuerpo $\K = (\enteros/7\enteros)$}:
        $$
          X^2 -2X + 10 =
          X^2 + 5X + 3
          \Entonces{raíces}
          r_{1,2} = \frac{2 \pm \sqrt{6}}{2} \notin (\enteros/7\enteros)
        $$
        Por lo tanto
        $X^2 -2X + 10$  también es irreducible en $(\enteros/7\enteros)$.
\end{enumerate}

\begin{aportes}
  \item \aporte{\dirRepo}{naD GarRaz \github}
  \item \aporte{https://github.com/ivdou}{Ivan Doumerc \github}
\end{aportes}
